\documentclass[12pt]{article}
\usepackage[utf8]{inputenc}
\usepackage[T1]{fontenc}
\usepackage{hyperref}
\title{Banam Bay Language Resources for K-12 Curriculum Development}
\author{Compiled by Claude Sonnet 5 (claude-sonnet-5) for Andr\'as Kornai}
\date{}
\begin{document}
\maketitle
\section{Introduction}
Banam Bay (ISO 639-3 vrt; Glottocode burm1263; also Vartavo, Burmbar) is an Oceanic language of south-east Malekula, Vanuatu, with fewer than about 900 speakers. It is poorly documented; the main effort found is a recent documentation project.
\section{Primary Educational Resources}
None found. The ELDP project is reported to aim at literacy materials, but none were seen in the search results.
\section{Linguistic and Reference Materials}
The ELDP-funded project (grant IGS0329) by Brittany Rae Hoback, Victoria University of Wellington, aims at a descriptive grammar and a standard orthography: \url{https://eldp.access.preservica.com/dk0549/}. Older sources have only recorded wordlists.
\section{Religious and Community Resources}
None found.
\section{Digital and Audio Resources}
The documentation collection is archived with ELDP/ELAR (link above). Endangered Languages Project entry: \url{https://www.endangeredlanguages.com/node/110139}
\section{Curriculum Development Context}
No evidence of school use was found. Any curriculum would depend on the outputs of the documentation project, particularly the orthography and grammar, and on community decisions.
\end{document}